% =====================================================================
%  5. Discussion and future work
% =====================================================================
\section{Discussion and Future Work}
\label{sec:discussion}

\subsection{Discussion}
\label{sec:disc}

All answers below come from the single test run on 2018 and agree with validation 2017.

\paragraph{Q1: feature engineering with a linear model.} Ridge with lags, calendar, holiday and
degree-day features removes about a third of the seasonal-naive error (880\,MW, 3.0\% of mean load) in
every month of the test year. This is by far the largest single gain in the study.

\paragraph{Q2: a non-linear model.} LightGBM removes a further 13\% of the Ridge error with actual
weather (231\,MW) and 10\% with lagged weather (168\,MW), winning in 10--11 of 12 months. The gain is
concentrated in cold and hot hours and in winter months, i.e.\ where the response to temperature
and its interaction with calendar matter most. Ridge stays better on bridge days, and LightGBM lost
September 2018, a month with unusual recorded load.

\paragraph{Q3: knowing the actual weather.} With perfect knowledge of the day's temperature,
LightGBM improves by 71\,MW (0.25\% of load); Ridge does not improve measurably. This is an upper
bound for the value of weather forecasts in this set-up, and it is smaller than the gain from
changing the model. It agrees with the weak within-month temperature effect found on the training
years (\cref{tab:s2}): once lags and calendar are present, temperature adds little extra information.
This does not mean temperature is unimportant; much of its effect is already carried by the season
and by recent load.

\paragraph{What this means for a public-data analyst.} Spend effort first on features (lags at the
right information time, calendar, holidays), then on a flexible model; a weather-forecast feed comes
last for a national day-ahead forecast of this kind. Holidays and bridge days remain the weakest
point of every model and are the most promising place for further work. The TSO forecast is still
far more accurate (268\,MW), so a public-data model is useful for understanding and analysis, not as
a replacement.

\subsection{Limitations}
\label{sec:limitations}

\begin{itemize}
  \item \textbf{No weather forecasts.} Only an optimistic (actual) and a pessimistic (lagged) bound
        can be measured; operational performance lies in between.
  \item \textbf{Few seasons.} Training covers two summers when tuning and three when testing. The
        June 2017 heat wave shows that unseen extremes cause large bias, especially for trees, which
        cannot extrapolate.
  \item \textbf{Weather quality.} A simple five-city mean, no population weights; the origin of the
        weather data is not documented.
  \item \textbf{Holidays.} Regional holidays are only captured as an unweighted share of regions.
  \item \textbf{One test year,} one refit per year (the model may age during 2018), heuristic
        thresholds ($\delta$, spike rule, grids).
  \item \textbf{Actual load} may have been revised after publication; the TSO column's information
        time is unknown.
  \item \textbf{Tuning budgets differ:} 94 Ridge fits versus 30 LightGBM trials whose early stopping
        uses the same validation year. The resulting optimism of M1 on validation is estimated at
        10--30\,MW, small next to the 175--249\,MW gap. The chosen \code{min\_data\_in\_leaf} (11--13)
        lies near the lower edge (10) of the search space. The number of trees is kept fixed when
        refitting on three years, which, if anything, slightly disadvantages M1.
  \item \textbf{Design details:} the spike rule looks at the next hour, so in principle the
        cleaning of hour 08 on $D-1$ uses hour 09 (no spike fell on hour 08). In the actual setting,
        the 24-hour mean temperature at early hours of $D$ also uses the rest of $D-1$. The holiday
        table misses five national holidays that fell on a Sunday in 2015--2017 (none in 2018).
        The eight test comparisons are not corrected for multiple testing.
  \item All relations are statistical associations, not causal effects.
\end{itemize}

\subsection{Future Work}
\label{sec:future}

Items are ordered by expected value for the questions of this report. Status is updated as work
progresses.

\begin{table}[H]
  \centering
  \small
  \caption{Planned extensions (not done yet unless marked).}
  \label{tab:future}
  \begin{tabularx}{\linewidth}{@{}l Y l@{}}
    \toprule
    ID & Extension & Status \\
    \midrule
    F1 & Monthly refit on the test year (rolling origin) instead of one refit per year. & \todo{} \\
    F2 & Population-weighted temperature, or degree days per city before averaging. & \todo{} \\
    F3 & Diebold--Mariano test as a cross-check of the block bootstrap. & \todo{} \\
    F4 & Ridge without load lags (2$\times$2: with/without lags $\times$ with/without temperature) to
         see how much temperature information the lags absorb. & \todo{} \\
    F5 & Error breakdown with and without special days, to compare with studies that exclude
         them~\cite{incremona2021italian,dudek2021ensembles}. & \todo{} \\
    F6 & Cubic temperature terms and temperature $\times$ month for Ridge (is Ridge limited by its
         functional form?). & \todo{} \\
    F7 & Humidity as an extra weather variable (actual setting only). & \todo{} \\
    F8 & Similar-day (pattern-based) baseline~\cite{dudek2021ensembles}. & \todo{} \\
    F9 & Sensitivity to an earlier issue time ($D-2$) for both load and weather. & \todo{} \\
    F10 & A deep model (MLP or LSTM) under the same protocol, if the scope allows. & \todo{} \\
    F11 & Probabilistic forecasts (quantile LightGBM) for operational use. & \todo{} \\
    F12 & Better bridge-day handling (e.g.\ a bridge $\times$ hour interaction, or using the previous
          year's bridge days), the weakest point of LightGBM in both years. & \todo{} \\
    F13 & Check the unusual load of September 2018 against the original ENTSO-E data (real events
          or date misalignment?). & \todo{} \\
    \bottomrule
  \end{tabularx}
\end{table}
